
\noindent Introduzido por \cite{vapnik1997support}, o método de Máquinas de Vetores de
Suporte (Support Vector Machines, SVM) constitui uma das abordagens mais influentes no
contexto do aprendizado estatístico. Sua principal característica reside na construção
de fronteiras de decisão potencialmente não lineares no espaço original de entrada, por
meio da indução de um hiperplano linear em um espaço de características de maior dimensão,
obtido a partir de uma transformação, possivelmente implícita, dos dados. Essa formulação
permite controlar a complexidade do modelo ao mesmo tempo em que busca maximizar a margem
de separação, o que confere ao método sólidas garantias teóricas de generalização.

Este capítulo apresenta o SVM nos contextos de classificação e regressão. Inicialmente,
é abordado o problema de classificação, cuja formulação estabelece os fundamentos teóricos
essenciais, tais como margem, vetores de suporte e dualidade, que são posteriormente
estendidos para os demais cenários considerados.

\section{Máquinas de Vetores de Suporte para Classificação (SVM)}

O objetivo do SVM para classificação é determinar um hiperplano que maximize a margem de
separação entre as classes. Formalmente, conforme apresentado por
\citeonline[p.~79]{mohri2018foundations}, considere um espaço de covariáveis
$\mathcal{X} \subseteq \mathbb{R}^d$ e um espaço de rótulos $\mathcal{Y} = \{-1, +1\}$.
Seja $\mathit{f} : \mathcal{X} \rightarrow \mathcal{Y}$ a função alvo desconhecida.

Dado um conjunto de hipóteses $\mathcal{H}$ de funções que mapeiam $\mathcal{X}$ em
$\mathcal{Y}$, o problema de classificação binária consiste em encontrar uma hipótese
$\mathit{h} \in \mathcal{H}$ que minimize o erro de generalização:
\[
\mathit{R}_{\mathcal{D}}(\mathit{h}) = \mathbb{P}_{\mathit{x} \sim \mathcal{D}} \left[\mathit{h}(\mathit{x}) \neq \mathit{f}(\mathit{x})\right],
\]
em que $\mathcal{D}$ é uma distribuição desconhecida sobre $\mathcal{X}$.

Para isso, o modelo recebe um conjunto de treinamento
\[
\mathit{S} = \left((\mathit{x}_1, \mathit{y}_1), \ldots, (\mathit{x}_m, \mathit{y}_m)\right) \in (\mathcal{X} \times \mathcal{Y})^m,
\]
composto por $\mathit{m}$ observações i.i.d. geradas segundo a distribuição $\mathcal{D}$,
tais que $\mathit{y}_i = \mathit{f}(\mathit{x}_i)$.

Um conjunto de hipóteses natural para esse problema é dado por classificadores lineares,
definidos por hiperplanos. Nesse caso,
\[
\mathcal{H} = \left\{\mathbf{x} \mapsto \sgn(\mathbf{w} \cdot \mathbf{x} + b) : \mathbf{w} \in \mathbb{R}^d, \; b \in \mathbb{R} \right\},
\]
em que $\mathbf{w} \cdot \mathbf{x} + b = 0$ representa um hiperplano em $\mathbb{R}^d$,
com $\mathbf{w} \in \mathbb{R}^d$ um vetor normal não nulo ao hiperplano e
$b \in \mathbb{R}$ um escalar. A regra de decisão associada é dada por
\[
\mathit{h}(\mathbf{x}) = \sgn(\mathbf{w} \cdot \mathbf{x} + b),
\]
conforme ilustrado na Figura \ref{fig:jabref}.

\begin{figure}[htb]
    \centering
    \includegraphics[scale=0.45]{hyperplane.png}
    \caption{Exemplo de hiperplanos de separação. O hiperplano à direita maximiza a margem entre as classes.}
    \label{fig:jabref}
\end{figure}

Em geral, podem existir infinitos hiperplanos que separam corretamente os dados. No entanto, o
SVM seleciona aquele que maximiza a margem geométrica entre as classes.

A margem geométrica $\rho_{\mathit{h}}(\mathbf{x})$ de um classificador $\mathit{h}$
em um ponto $\mathbf{x}$ é definida como a distância Euclidiana desse ponto ao hiperplano:
\[
\rho_{\mathit{h}}(\mathbf{x}) = \frac{|\mathbf{w} \cdot \mathbf{x} + b|}{\|\mathbf{w}\|_2}.
\]

Para uma amostra $\mathit{S} = (\mathbf{x}_1, \ldots, \mathbf{x}_m)$, a margem geométrica do
classificador é definida como o menor valor entre as margens individuais:
\[
\rho_{\mathit{h}} = \min_{i \in [m]} \rho_{\mathit{h}}(\mathbf{x}_i),
\]
o que corresponde à distância do hiperplano aos pontos mais próximos da amostra.

Assim, o problema de maximização da margem pode ser formulado como
\[
\rho = \max_{\mathbf{w}, b} \min_{i \in [m]} \frac{\mathit{y}_i (\mathbf{w} \cdot \mathbf{x}_i + b)}{\|\mathbf{w}\|},
\]
sob a condição de correta classificação dos dados.

Observe que essa expressão é invariante à multiplicação de $(\mathbf{w}, b)$ por um escalar
positivo. Dessa forma, pode-se impor a normalização
\[
\min_{i \in [m]} \mathit{y}_i (\mathbf{w} \cdot \mathbf{x}_i + b) = 1,
\]
o que conduz à formulação equivalente
\[
\rho = \max_{\mathbf{w}, b} \frac{1}{\|\mathbf{w}\|} \quad \text{sujeito a} \quad \mathit{y}_i (\mathbf{w} \cdot \mathbf{x}_i + b) \geq 1, \; \forall i \in [m].
\]

Como maximizar $\frac{1}{\|\mathbf{w}\|}$ é equivalente a minimizar $\frac{1}{2}\|\mathbf{w}\|^2$,
obtém-se o seguinte problema de otimização convexa:
\begin{equation}
\begin{aligned}
\min_{\mathbf{w}, b} \quad & \frac{1}{2} \|\mathbf{w}\|^2 \\
\text{sujeito a:} \quad & \mathit{y}_i (\mathbf{w} \cdot \mathbf{x}_i + b) \geq 1, \quad \forall i \in [m].
\end{aligned}
\label{eq:svm-problem}
\end{equation}

A função objetivo $F(\mathbf{w}) = \frac{1}{2} \|\mathbf{w}\|^2$ é estritamente convexa,
uma vez que sua matriz Hessiana é a identidade, que é definida positiva. As restrições
são funções afins e, portanto, convexas. Assim, o problema definido em \ref{eq:svm-problem}
é um problema de otimização convexa que admite solução ótima única.

Essa propriedade constitui uma vantagem fundamental do SVM, pois garante que a solução
encontrada corresponde ao ótimo global e pode ser obtida de forma eficiente por métodos de
otimização convexa.

A formulação apresentada pode ser reescrita em termos de variáveis Lagrangeanas, conduzindo
ao problema dual. No entanto, essa abordagem será desenvolvida na próxima seção, no contexto
de dados não linearmente separáveis, que é mais representativo de aplicações reais.

\subsection{Caso não separável}

Na prática, é comum que os dados não sejam linearmente separáveis, isto é, exista ao menos
um índice $i \in \left[\mathit{m}\right]$ tal que $\mathit{y}_i \left(\mathbf{w} \cdot \mathbf{x}_i + b\right) < 1$,
como ilustrado na Figura \ref{fig:non-separable}. Nesse cenário, introduzem-se variáveis
de folga $\xi_i \geq 0$, que permitem a violação controlada das restrições de margem,
tornando o problema mais flexível e robusto a ruídos e outliers.

Dessa forma, as restrições de separação passam a ser escritas como
$$
\mathit{y}_i \left(\mathbf{w} \cdot \mathbf{x}_i + b\right) \geq 1 - \xi_i, \quad \forall i \in \left[\mathit{m}\right].
$$

\begin{figure}[htb]
    \centering
    \includegraphics[scale=0.45]{hyperplan_non_separable.png}
\caption{Hiperplano de separação no caso não separável. As variáveis de folga $\xi_i$ permitem a violação das restrições de margem.}
\label{fig:non-separable}
\end{figure}

A variável $\xi_i$ quantifica o grau de violação da restrição associada à observação
$\mathbf{x}_i$. Em particular, quando $0 < \xi_i < 1$, o ponto é corretamente classificado,
mas encontra-se dentro da margem; quando $\xi_i \geq 1$, ocorre erro de classificação.
Assim, as variáveis de folga permitem distinguir diferentes níveis de violação das condições
ideais de separação.

Nesse contexto, surge um compromisso entre maximizar a margem geométrica e penalizar
violações, o que leva à seguinte formulação de otimização:

\begin{equation}
\begin{aligned}
\min_{\mathbf{w}, b, \bm{\xi}} \quad & \frac{1}{2} \|\mathbf{w}\|^2 + C \sum_{i=1}^{m} \xi_i \\
\text{sujeito a:} \quad & \mathit{y}_i \left(\mathbf{w} \cdot \mathbf{x}_i + b\right) \geq 1 - \xi_i, \quad \forall i \in \left[\mathit{m}\right], \\
& \xi_i \geq 0, \quad \forall i \in \left[\mathit{m}\right],
\end{aligned}
\label{eq:svm-non-separable}
\end{equation}

em que $C > 0$ é um parâmetro de regularização que controla o equilíbrio entre a maximização
da margem e a penalização das violações.

O problema em \eqref{eq:svm-non-separable} é um programa quadrático convexo, uma vez que a
função objetivo é estritamente convexa e as restrições são afins. Para derivar sua forma
dual, introduzem-se multiplicadores de Lagrange $\alpha_i \geq 0$ associados às restrições
de margem e $\mu_i \geq 0$ associados às restrições $\xi_i \geq 0$. A função Lagrangiana é
então dada por:

\begin{equation}
\begin{aligned}
L(\mathbf{w}, b, \bm{\xi}, \bm{\alpha}, \bm{\mu}) =
\frac{1}{2} \|\mathbf{w}\|^2
+ C \sum_{i=1}^{m} \xi_i
- \sum_{i=1}^{m} \alpha_i \left[\mathit{y}_i (\mathbf{w} \cdot \mathbf{x}_i + b) - 1 + \xi_i \right]
- \sum_{i=1}^{m} \mu_i \xi_i.
\end{aligned}
\end{equation}

As condições de Karush-Kuhn-Tucker (KKT) são obtidas impondo estacionariedade em relação
às variáveis primais e complementaridade:

\begin{align}
\nabla_{\mathbf{w}} L &= \mathbf{w} - \sum_{i=1}^{m} \alpha_i \mathit{y}_i \mathbf{x}_i = 0
&&\Longrightarrow\quad
\mathbf{w} = \sum_{i=1}^{m} \alpha_i \mathit{y}_i \mathbf{x}_i,
\label{eq:kkt_w_ns} \\[6pt]
\nabla_{b} L &= - \sum_{i=1}^{m} \alpha_i \mathit{y}_i = 0
&&\Longrightarrow\quad
\sum_{i=1}^{m} \alpha_i \mathit{y}_i = 0,
\label{eq:kkt_b_ns} \\[6pt]
\nabla_{\xi_i} L &= C - \alpha_i - \mu_i = 0
&&\Longrightarrow\quad
0 \leq \alpha_i \leq C,
\label{eq:kkt_xi_ns} \\[6pt]
&\alpha_i \left[\mathit{y}_i (\mathbf{w} \cdot \mathbf{x}_i + b) - 1 + \xi_i \right] = 0,
\label{eq:kkt_comp1_ns} \\[6pt]
&\mu_i \xi_i = 0.
\label{eq:kkt_comp2_ns}
\end{align}

Da condição \eqref{eq:kkt_w_ns}, observa-se que $\mathbf{w}$ pertence ao subespaço gerado
pelos dados de treinamento. Além disso, apenas observações com $\alpha_i > 0$ contribuem
para a solução, sendo denominadas vetores de suporte. Esses pontos são precisamente aqueles
que estão sobre ou dentro da margem, incluindo possíveis violações.

Substituindo as condições de estacionariedade na Lagrangiana, obtém-se o problema dual:

\begin{equation}
\begin{aligned}
\max_{\bm{\alpha}} \quad & \sum_{i=1}^{\mathit{m}} \alpha_i - \frac{1}{2} \sum_{i,j=1}^{\mathit{m}} \alpha_i \alpha_j \mathit{y}_i \mathit{y}_j \left(\mathbf{x}_i \cdot \mathbf{x}_j\right) \\
\text{sujeito a:} \quad & \sum_{i=1}^{\mathit{m}} \alpha_i \mathit{y}_i = 0, \\
& 0 \leq \alpha_i \leq C, \quad \forall i \in \left[\mathit{m}\right].
\end{aligned}
\end{equation}

A função objetivo dual é côncava, pois sua Hessiana é dada por $-\mathbf{Q}$, com
$\mathbf{Q}$ sendo semidefinida positiva. Assim, o problema dual é um problema de otimização
convexa bem posto, com solução ótima global.

A hipótese final pode ser expressa como:
\begin{equation}
\mathit{h}(\mathbf{x}) = \sgn \left( \sum_{i=1}^{\mathit{m}} \alpha_i \mathit{y}_i (\mathbf{x}_i \cdot \mathbf{x}) + b \right).
\end{equation}

Essa formulação evidencia que a solução depende exclusivamente de produtos internos entre
observações, o que possibilita a generalização para espaços de alta dimensão por meio de
funções de kernel.

Por fim, a relação entre a norma de $\mathbf{w}$ e os multiplicadores de Lagrange permite
interpretar a margem geométrica em termos da norma $\ell_1$ de $\bm{\alpha}$, reforçando
que a solução do SVM é essencialmente determinada pelos vetores de suporte.

Com base nesses fundamentos, a próxima seção apresenta a extensão natural do SVM para
problemas de regressão, na qual a noção de margem é adaptada para lidar com variáveis
resposta contínuas.

\section{Máquinas de Vetores de Suporte para Regressão (SVR)}
\label{sec:svr}

No contexto de regressão, o objetivo das Máquinas de Vetores de Suporte consiste em
aproximar uma função a partir dos dados de treinamento, controlando simultaneamente a
complexidade do modelo e o erro de ajuste. Diferentemente do caso de classificação,
em que se busca uma fronteira de decisão, no SVR procura-se ajustar uma função tal que
os desvios em relação aos valores observados sejam, sempre que possível, inferiores a um
limiar previamente especificado $\epsilon > 0$.

Geometricamente, essa ideia pode ser interpretada como o ajuste de um tubo de largura
$2\epsilon$ em torno da função estimada, conforme ilustrado na Figura \ref{fig:svm-regression}.
Observações cujos valores preditos se encontram no interior desse tubo não são penalizadas,
enquanto aquelas situadas fora do tubo incorrem em penalizações proporcionais à magnitude
da violação.

\begin{figure}[htb]
    \centering
    \includegraphics[scale=0.45]{svm_regression.png}
\caption{SVM para regressão. O objetivo é ajustar um tubo de largura $2\epsilon$ de modo que a maioria das observações esteja contida em seu interior.}
\label{fig:svm-regression}
\end{figure}

Formalmente, considere o conjunto de hipóteses
$$
\mathcal{H} = \{x \mapsto \mathbf{w} \cdot \Phi(x) + b : \mathbf{w} \in \mathbb{R}^N, b \in \mathbb{R}\},
$$
em que $\Phi$ é uma aplicação de mapeamento associada a um kernel definido positivo
(PDS) $K$. O problema de otimização do SVR pode ser formulado como:

\begin{equation}
\min_{\mathbf{w}, b} \quad \frac{1}{2} \|\mathbf{w}\|^2 + C \sum_{i=1}^{\mathit{m}} |\mathit{y}_i - (\mathbf{w} \cdot \Phi(x_i) + b)|_\epsilon,
\end{equation}

em que $|\cdot|_{\epsilon}$ denota a perda $\epsilon$-insensível, definida por:
\begin{equation}
|\mathit{y} - \mathit{y}'|_\epsilon = \max\{0, |\mathit{y} - \mathit{y}'| - \epsilon\}.
\end{equation}

Essa função de perda ignora erros de pequena magnitude, promovendo soluções esparsas,
uma vez que apenas observações fora do tubo $\epsilon$ contribuem para a função objetivo.

Para obter uma formulação tratável, introduzem-se variáveis de folga
$\xi_i \geq 0$ e $\xi_i' \geq 0$, que capturam, respectivamente, desvios acima e abaixo
do tubo. O problema pode então ser reescrito como:

\begin{equation}
\begin{aligned}
\min_{\mathbf{w}, b, \bm{\xi}, \bm{\xi}'} \quad
& \frac{1}{2} \|\mathbf{w}\|^2 + C \sum_{i=1}^{\mathit{m}} (\xi_i + \xi_i') \\
\text{sujeito a:} \quad
& \mathbf{w} \cdot \Phi(x_i) + b - \mathit{y}_i \leq \epsilon + \xi_i, \\
& \mathit{y}_i - \mathbf{w} \cdot \Phi(x_i) - b \leq \epsilon + \xi_i', \\
& \xi_i \geq 0, \quad \xi_i' \geq 0, \quad \forall i \in \left[\mathit{m}\right].
\end{aligned}
\end{equation}

Esse problema é um programa quadrático convexo com restrições afins. Introduzindo
multiplicadores de Lagrange e aplicando as condições de Karush-Kuhn-Tucker, obtém-se
a seguinte formulação dual:

\begin{equation}
\begin{aligned}
\max_{\bm{\alpha}, \bm{\alpha}'} \quad
& \varepsilon \sum_{i=1}^{m} (\alpha_i + \alpha_i')
- \sum_{i=1}^{m} y_i (\alpha_i' - \alpha_i) \\
& - \frac{1}{2} \sum_{i,j=1}^{m} (\alpha_i' - \alpha_i)(\alpha_j' - \alpha_j) K_{ij} \\
\text{sujeito a:} \quad
& \sum_{i=1}^{m} (\alpha_i' - \alpha_i) = 0, \\
& 0 \leq \alpha_i \leq C, \quad 0 \leq \alpha_i' \leq C, \quad \forall i \in \left[\mathit{m}\right],
\end{aligned}
\end{equation}

em que $K_{ij} = K(x_i, x_j) = \Phi(x_i) \cdot \Phi(x_j)$.

A hipótese resultante pode ser expressa como:
\begin{equation}
h(x) = \sum_{i=1}^{m} (\alpha_i' - \alpha_i) K(x_i, x) + b.
\end{equation}

O termo $b$ pode ser estimado a partir de condições de complementaridade,
utilizando pontos para os quais $0 < \alpha_i < C$ ou $0 < \alpha_i' < C$.

As condições de complementaridade implicam que, para cada $i$, no máximo um dos
multiplicadores $\alpha_i$ ou $\alpha_i'$ é estritamente positivo, refletindo o fato
de que uma observação não pode simultaneamente violar o tubo em ambas as direções.
Além disso, para pontos situados no interior do tubo $\epsilon$, tem-se
$\alpha_i = \alpha_i' = 0$, de modo que tais observações não contribuem para a solução.

Consequentemente, apenas os pontos que violam a margem $\epsilon$ influenciam
diretamente a hipótese final, caracterizando-se como vetores de suporte. Quando uma
fração significativa dos dados encontra-se no interior do tubo, a solução tende a ser
esparsa, o que é uma propriedade desejável do ponto de vista computacional e estatístico.

O parâmetro $\epsilon$ desempenha um papel central nesse contexto, estabelecendo um
compromisso entre esparsidade e qualidade de ajuste. Valores maiores de $\epsilon$
resultam em soluções mais esparsas, ao custo de uma possível perda de precisão, enquanto
valores menores aumentam a sensibilidade do modelo aos dados.

O SVR apresenta diversas propriedades favoráveis. Em particular, está fundamentado em
princípios teóricos sólidos, admite soluções esparsas e permite a incorporação natural
de kernels, possibilitando a modelagem de relações não lineares em espaços de alta
dimensão. Ademais, o método apresenta propriedades de estabilidade que garantem boa
capacidade de generalização, conforme discutido em \citeonline[p.~339--340]{mohri2018foundations}.

Por outro lado, o desempenho do SVR depende da escolha adequada dos hiperparâmetros
$C$ e $\epsilon$, cuja seleção usualmente requer procedimentos de validação cruzada,
potencialmente custosos em termos computacionais. Além disso, a resolução do problema
de otimização pode se tornar onerosa para conjuntos de dados de grande escala.
